\section{Discussion}
\label{sec:discussion}

\paragraph{Sumstats-only validation is rigorous but lower-power.}
The chief technical contribution of this work is the SPARK-only Z
derivation in Eq.~\ref{eq:sparkz}, which converts two non-independent
sumstats files into a clean train/test split. Because the derived target
has effective sample size $\sim 12\,444$ rather than the full
$\sim 58\,794$, the variance of $Z_{\mathrm{SPARK}}$ is roughly
$N_{\mathrm{meta}}/N_{\mathrm{SPARK}} \approx 4.7\times$ that of the
meta-analysis Z. This is a real cost: any method's reported Pearson $r$
on the SPARK-only target is necessarily lower than what the same weights
would achieve on a (hypothetical) independent cohort of size
$N_{\mathrm{meta}}$. Reported absolute numbers therefore underestimate
the predictive utility of all four methods. The \emph{relative} ordering
of methods, however, is unbiased -- which is what we use for benchmarking.

\paragraph{Why the ML re-weighters win even without annotations.}
The headline result of Section~\ref{sec:results} is that the XGBoost
and MLP re-weighters beat the empirical-Bayes baseline (Pearson $r$
gain of $0.017$ and $0.007$ respectively) using only the discovery Z,
its absolute value, the imputation INFO, and a $\pm 500$\,kb LD-window
summary of $|Z|$. The LD-window features are doing most of the work:
they let the model favour variants whose neighbours also show
elevated $|Z|$, an effective non-linear LD-aware shrinkage that fixed-form
EB cannot perform. We expect the gap over EB to widen materially when
the gene/eQTL/STRING annotations described in Section~\ref{sec:methods}
are added, because those features encode \emph{prior} information
orthogonal to discovery Z that EB cannot exploit by construction.

\paragraph{Where the ML methods should help and where they should not.}
With annotations available, the feature-based re-weighter and the
gene-network GNN both add a prior that variants near brain-relevant,
conserved, or SFARI-network-central genes carry larger true effects.
This prior should help when (i) the discovery Z is small but the
annotation evidence is strong (sub-significant variants in important
genes) and (ii) the discovery Z is moderate but the gene-network
context confirms a wider biological pattern. The same prior should
help \emph{less} on variants whose nearest gene is poorly annotated or
absent from STRING.

\paragraph{Limitations.}
\begin{itemize}[nosep,leftmargin=1.2em]
  \item \emph{Modest absolute Pearson $r$.} Reported $r$ values
        ($0.04$--$0.06$) reflect the high per-SNP measurement noise
        inherent in the derived SPARK-only Z at
        $N_{\mathrm{SPARK}} = 12\,444$. Variance in
        $Z_{\mathrm{SPARK}}$ is roughly $4.7\times$ that of the meta
        Z, so the same weights would achieve substantially higher
        concordance on a hypothetical independent cohort of size
        $N_{\mathrm{meta}}$. The relative ordering of methods is
        unaffected by this scaling.
  \item \emph{Annotation files were not available locally for this
        run.} The gene-network GNN therefore reduced to the discovery
        baseline because no SNP could be mapped to a gene, and the ML
        re-weighters relied only on the four "self'' features listed
        above. The ranking already favours the ML methods; the gain
        will widen once the annotations are downloaded.
  \item \emph{No individual-level evaluation.} The headline metrics are
        SNP-level concordance, not per-individual case/control
        prediction. Translating these into liability-scale $R^2$ on a
        real cohort requires individual-level data (e.g., SPARK
        genotypes via SFARI Base, dbGaP, or UK Biobank).
  \item \emph{Single ancestry.} Both releases are predominantly
        European-ancestry; cross-ancestry transferability is not
        addressed here.
  \item \emph{Dependence on functional annotations.} Without GTEx,
        ENCODE, STRING, and the SFARI Gene CSV, the ML re-weighters
        collapse onto the LD-window features and an in-gene boolean,
        and most of the additive value of the approach is lost.
  \item \emph{Linear projection at the gene/SNP boundary.} The GNN's
        re-projection to SNPs uses a fixed $w = 0.5$ convex combination;
        a learned per-gene weighting could improve calibration further.
\end{itemize}

\paragraph{Future work.}
Beyond the obvious extension to individual-level genotypes once a SPARK
or dbGaP DUA is granted, the most natural directions are: (i)
multi-trait modelling that jointly fits ASD with co-heritable conditions
(ADHD, schizophrenia, intelligence) under a shared annotation backbone;
(ii) integration of rare-variant gene burden scores at the gene-graph
nodes, formally combining the common- and rare-variant axes that
Dong et al.~\cite{dong2023} treat in parallel; and (iii) a
better-calibrated SNP-level re-projection that adapts $w$ to local LD
density and the GNN's own predictive uncertainty.
